\documentclass[11pt]{article}
% comp_certcup_en standalone skeleton (archive comp_certcup_en: template_cls =
% article; rules: English only, Summary on a standalone page, structure follows
% the US-style outline, pdflatex, <= 25 pages advisory). This skeleton belongs
% to this contest only: it carries no other contest's cover, control number, or
% declaration. Standalone since 2026-09-27 (dispatcher-authorized cleanup of a
% reused-skeleton contamination; see C-window REPORT_ROUND6).
\usepackage[a4paper,margin=2.5cm]{geometry}
\usepackage{newtxtext}

% === Extra packages ===
\usepackage{algorithm2e}
\usepackage{float}
\usepackage{longtable}
\usepackage{tabularx}
\usepackage[section]{placeins}
\usepackage{needspace}

% === TikZ ===
\usepackage{tikz}
\usetikzlibrary{arrows.meta,positioning,shapes.geometric,fit,calc}

\begin{document}

% === Summary page (archive rule: Summary must be a standalone page) ===
\begin{center}
    {\LARGE\bfseries [Paper Title]}
\end{center}

\vspace{1.5em}

\begin{center}
    {\large\bfseries Summary}
\end{center}

\noindent [Summary: 300--400 words on one page; one paragraph per sub-problem with a named method and verified numerical result; use \textbf{} for 1--3 concise method/result anchors per sub-problem]

\vspace{1.5em}
\noindent\textbf{Keywords:} [Keyword1]; [Keyword2]; [Keyword3]

\newpage
% ⛔ No table-of-contents page: contest papers do not carry one.
%   The \newpage above keeps the Summary page standalone before the body.

% === Body ===
\input{sections/1_introduction}
\input{sections/2_assumptions}
\input{sections/3_model_design}
\input{sections/4_solution}
\input{sections/5_sensitivity}
\input{sections/6_strengths_weaknesses}
\input{sections/7_conclusions}

% === References ===
\begin{thebibliography}{99}
% \bibitem{ref1} Author. Title[J]. Journal, Year.
\end{thebibliography}

% === Appendix ===
\appendix
\input{sections/A_code}

\end{document}
